\documentclass[english,twoside, openright, hidelinks, 11 pt, class=report,crop=false]{standalone}
\input{../../preamb}
\input{../bm}

\begin{document}
	
\op{opgalg}
Vis at $n^2+n$ er et partall.	
	
\op{opggeomidtparl}
$E$, $F$, $G$ og $H$ er de respektive midtpunktene til sidene til $\square ABCD$.
\abc{
\item Vis at $\square EFGH$ er et parallellogram.
\item Vis at $A_{\square EFGH}=\frac{1}{2}A_{\square ABCD}$
}
\fig{opggeomidtparl}
\mers{Sammenhengen fra oppgave a) kalles \outl{Varignons teorem}.}

\newpage
\sv
\fig{opggeomidtparl_lf}
\abc{
\item $\triangle ACD\sim\triangle HGD$ fordi $\angle ADC=\angle HDG$ og $\frac{AD}{HD}=\frac{DC}{DG}=2$. Tilsvarende er $\triangle ABC\sim \triangle EBF$. Dermed er både $HG$ og $EF$ parallelle med $AC$. Tilsvarende kan vi vise at $HE$ og $GF$ er parallelle med $BD$. Dermed er $\square EFGH$ et parallellogram.
\item Vi har at
\begin{align}
	2A_{\square ABCD} &= A_{\triangle ACD}+A_{\triangle ABC}+A_{\triangle ABD}+A_{\triangle BCD} \label{opggeomidtparlabcd}
\end{align}
og at 
\begin{align}
	A_{\square EFGH}&=A_{\square ABCD}-A_{\triangle HGD}-A_{\triangle EBF}-A_{\triangle AEH}-A_{\triangle FCH} \label{opggeomidtparlefgh}
\end{align}
Av formlikheten vi viste i oppgave a), vet vi at
\[ A_{\triangle HGD}=\left(\frac{1}{2}\right)^2 A_{\triangle ACD}=\frac{1}{4}A_{\triangle ACD}\]
Tilsvarende sammenheng har vi mellom de andre trekantene fra \eqref{opggeomidtparlabcd} og trekantene fra \eqref{opggeomidtparlefgh}, og dermed er
\begin{align*}
	A_{\square EFGH}&=A_{\square ABCD}- \frac{1}{4}\cdot 2A_{\square ABCD}=\frac{1}{2}A_{\square ABCD}
\end{align*}
}

\newpage
\fig{opggeoab_los}
Da de er periferivinkler som spenner over buer med lik lengde, er $\angle ADB=\angle BDC=\angle ACB=60^\circ$. Tilsvarende er $\angle CAD=\angle CBD=u$. Vi plasserer $F$ på forlengelsen av $AD$ slik at $\angle DFC=60^\circ$. Da er $\triangle CFD$ likesidet. Videre plasserer vi $E$ slik at $\triangle BCE\cong\triangle CAD$. Av vinkelsummen til $\triangle CAD$, er $\angle DCA=60^\circ-u$, og følgelig er $\angle FCE=180^\circ$. Videre har vi at $\angle EBD=\angle EBC+\angle CBD=60^\circ-u+u=60^\circ$. Følgelig er $\square FDEB$ et parallellogram ($\angle CEB=\angle BDF=120^\circ$), og da er 
\[ BD=EF=FC+CE=CD+DA \]

\newpage
\fig{opggeoab_los}
Since they are peripheral angles spanning arcs of equal length, $\angle ADB = \angle BDC = \angle ACB = 60^\circ$. Similarly, $\angle CAD=\angle CBD=u$. We place $F$ on the extension of $AD$ such that $\angle DFC = 60^\circ$. Then, $\triangle CFD$ is equilateral. Furthermore, we place $E$ such that $\triangle BCE \cong \triangle CAD$. From the angle sum of $\triangle CAD$, $\angle DCA = 60^\circ - u$, and consequently, $\angle FCE = 180^\circ$. Additionally, we have $\angle EBD = \angle EBC + \angle CBD = 60^\circ - u + u = 60^\circ$. Consequently, quadrilateral $FDEB$ is a parallelogram ($\angle CEB = \angle BDF = 120^\circ$), and thus
\[ BD = EF = FC + CE = CD + DA. \]

\newpage
\grubop{opggeotricharxandy}
Gitt en trekant $\triangle ABC$, hvor horisontalavstanden mellom $A$ og $B$ er $x$. Videre er $y=CD$ vertikalavstanden mellom $C$ og $AB$, eller mellom $C$ og forlengelsen av $AB$.\os

Vis at
\[ A_{\triangle ABC}=\frac{xy}{2} \]
\begin{figure}
	\centering
	\subfloat[]{\includegraphics{\figp{opggeotricharxandy_a}}} \qquad
	\subfloat[]{\includegraphics{\figp{opggeotricharxandy_b}}}
\end{figure}

\sv
\textbf{\boldmath Tilfellet hvor $D$ ligger på $AB$}
\fig{opggeotricharxandy_alf}
Vi danner parallellogrammet $\square ABEF$, hvor $AF\parallel EB \parallel CD$. Med $AB$ som grunnlinje, har $\square ABEF$ og $\triangle ABC$ lik høyde, og dermed er
\[ A_{\triangle ABC}=\frac{1}{2}A_{\square ABEF}= \frac{1}{2}AF\cdot GE=\frac{yx}{2} \]
\newpage
\textbf{\boldmath Tilfellet hvor $D$ ligger på forlengelsen av $AB$}
\fig{opggeotricharxandy_blf}
$C'$ er $C$ parallellforskjøvet langs en linje som er parallell med $AB$, slik at $D'$ ligger på $AB$. Da vet vi at
\[ A_{\triangle ABC'}=\frac{xy}{2} \]
Med $AB$ som grunnlinje, har $\triangle ABC'$ og $\triangle ABC$ lik høyde, og dermed likt areal. Følgelig er
\[ A_{\triangle ABC}=A_{\triangle ABC'}=\frac{xy}{2} \]
\newpage
Given a triangle $\triangle ABC$, where the horizontal distance between $A$ and $B$ is $x$. Furthermore, let $y = CD$ denote the vertical distance from $C$ to the line segment $AB$, or to the extension of $AB$.\os

Show that
\[
A_{\triangle ABC}=\frac{xy}{2}.
\]

\begin{figure}
	\centering
	\subfloat[]{\includegraphics{\figp{opggeotricharxandy_a}}} \qquad
	\subfloat[]{\includegraphics{\figp{opggeotricharxandy_b}}}
\end{figure}

\newpage
Sirkelen har radius $ r $.	Vis at $ 4BC\cdot r=\sqrt{2}BE\cdot CF $	
\fig{opggeo4zr}
Let $ r $ denote the radius of the circle.	Show that
\[ 4BC\cdot r=\sqrt{2}BE\cdot CF \]	
\fig{opggeo4zr}

\newpage
\abc{
\item Given the functions $f(x)=4x$ and $g(x)=-\frac{1}{4}x$. Show that the graph of $f$ and the graph of $g$ are perpendicular.
\item Given the functions $f(x)=ax+b$ and $g(x)=-\frac{1}{a}x+c$. Show that the graph of $f$ and the graph of $g$ are perpendicular.
}


\newpage
\abc{
\item \textbf{Method 1}\\
The graphs of $f$ and $g$ intersect when $x=0$. We set\\ $A=(0, 0)$, $B=(1, 0)$, and $C=(1, f(1))=(1, 4)$. Furthermore, we set $D=(-4, 0)$ and $E=(-4, g(-4))=(-4, 1)$. Since $DE=AB$ and $DA=BC$, and $\angle ADE=\angle CBA$, the triangles $\triangle EDA \sim \triangle ABC$.
\fig{opggeoperlina1_lf}
\newpage
\textbf{Method 2} \\
The graphs of $f$ and $g$ clearly intersect when $x=0$. We set\\ $A=(0, 0)$, $B=(1, 0)$, and $C=(1, f(1))=(1, 4)$. Let $E$ be the point on the graph of $g$ with the same $y$-value as $C$. At $E$,
\alg{
	g(x)&=4 \\
	-\frac{1}{4}x &= 4 \\
	x &= -16
}
Thus, $E=(-16, 4)$. We set $D=(-16, 0)$. By the Pythagorean theorem on $\triangle AED$, we have
\[ EA^2 = ED^2 + AD^2 = 4^2 + 16^2 = 16 + 16^2 = 16(1 + 16) = 16 \cdot 17 \]
By the Pythagorean theorem on $\triangle BAC$, we have
\[ AC^2 = AB^2 + BC^2 = 1^2 + 4^2 = 17 \]
If the graphs of $f$ and $g$ are perpendicular to each other, then $\angle CAE = 90^\circ$, and $\triangle AEC$ must satisfy the Pythagorean theorem. We know that $EC = BD = 17$, and thus we must require that
\[ EC^2 = AC^2 + EA^2 = 17^2 \]
We have
\[ AC^2 + EA^2 = 17 + 16 \cdot 17 = 17(1 + 16) = 17^2 \]
Thus, $\triangle AEC$ is a right triangle.
\fig{opggeoperlina2_lf}
\newpage
\item We assume $a > 1$. If $f = g$, we have
\begin{align*}
	ax + b &= -\frac{1}{a}x + c \br
	x &= \frac{a(b - c)}{a^2 + 1}
\end{align*}
This equation is guaranteed to have a solution, and thus we know that $f$ and $g$ intersect at a point $A$. We place $B$ and $D$ on the same horizontal line as $A$ such that $AD = BD = 1$. Furthermore, we place $E$ and $C$ on the graphs of $g$ and $f$, respectively, such that $DA \perp ED$ and $AB \perp BC$. Then, $DE = \frac{1}{a}$ and $BC = a$. We know that in $\triangle ADE$, the longest leg is $AD$ and the shortest leg is $DE$, and in $\triangle ABC$, the longest leg is $BC$ and the shortest leg is $AB$. Furthermore,
\[ \frac{DE}{AD} = \frac{AB}{BC} = \frac{1}{a} \]
Since $\angle ADE = \angle CBA$, it follows that $\triangle ADE \sim \triangle CBA$, and thus
\begin{align*}
	\angle CAE &= 180^\circ - \angle EAD - \angle BAC \\
	&= 180^\circ - \angle EAD - \left(90^\circ - \angle EAD\right) \\
	&= 90^\circ
\end{align*}
The explanations for the cases where $0 < a < 1$ or $a < 0$ are almost identical to the one we have given for $a > 1$.
\fig{opggeoperlinb_lf}
}

\grubop{opgbrokgrub}
For å finne hvilke av to brøker som er størst, kan vi utvide brøkene slik at de har lik nevner. Kan vi også avgjøre hvilken som er størst ved å utvide brøkene slik at de har lik teller? Grunngi svaret ditt.

\grubop{opggeofindh}
Gitt tre punkt. Finn avstanden fra det éne punktet til de to andre.

\grubop{opggeotrivek}
For trekant utspent av vektorene $[a, b]$ og $[c, d]$. Vis at $T=ad-bc$.

\grubop{opggeodistpoint}
Finn lengden til det grønne linjestykket.
\fig{opggeodistpoint}

\newpage
\fig{opggeodistpoint_lf}
Vi plasserer linjestykkene i et kvadratisk rutenett hvor rutene har sidelengde 1.
Vi har at
\begin{align*}
	A_{\triangle ABE}&=\frac{1}{2}AB\cdot AE=\frac{1}{2}\cdot12\cdot5 = 30 \vn
	A_{\triangle EDF}&=\frac{1}{2}ED\cdot DF=\frac{1}{2}\cdot1\cdot6 = 3 \vn
	A_{\triangle FCB}&=\frac{1}{2}FC\cdot BC=\frac{1}{2}\cdot6\cdot6=18\vn
	A_{\square ABCD}&=AB\cdot BC=12\cdot6=72
\end{align*}
Dermed er 
\begin{align*}
	A_{\triangle EBF}&=A_{\square ABCD}-A_{\triangle ABE}-A_{\triangle EDF}-A_{\triangle FCB} \\
	&= 72-30-3-18 \\
	&=21
\end{align*}
Av \pyt\ på $\triangle ABE$ har vi at
\begin{align*}
	BE^2 &= AE^2 + AB^2 \\
	&= 5^2+12^2 \\
	&= 169 \\
	BE &= 13
\end{align*}
Videre har vi at
\begin{align*}
	BE\cdot DF&=2A_{\triangle EBF} \\
	DF &= \frac{2A_{\triangle EBF}}{BE} \br
	DF &= \frac{2\cdot 21}{13} \br
	&= \frac{42}{13}
\end{align*}
\newpage

%alg
\grubeksoprecc{gv24d1opg7}{GV24}
Det røde kvadratet og den grønne sekskanten danner et kvadrat.
Vis at arealet til det grønne området er $a^2+2ab$.
\fig{gv24d1opg7}

%funk
\grubeksoprecc{gv22d2opg4}{GV22}
I et mønster er antall ruter $f(n)$ i figur nummer $n$ gitt ved
\[ f(n)=n^2+2 \]
Tegn det som kan være de tre første figurene i dette mønsterert.

%alg
\grubeksoprecc{gv22d2opg2}{GV22}
Figuren viser et stort kvadrat som er bygd opp av fire kongruente rektangler og et lite kvadrat.
Omkretsen til hvert rektangel er 30. \os

Argumenter for at omkretsen til det store kvadratet er 60 cm.
\fig{gv22d2opg2}

%alg
\grubeksoprecc{gv22d2}{GV22}
Vi har at
\[ 3^2=(4\cdot2)+1\quad,\quad 5^2=(6\cdot4)+1\quad,\quad 9^2=(10\cdot8)+1	 \]
Skriv denne sammenhengen uttrykt ved et heltall $n$, og vis at den gjelder for alle $n>1$.

\end{document}

